\scanpage{133}
3). The proof of formula (6) is the most difficult.
We begin by giving a
direct and very simple proof for the case where
\deleted{$\unsure{p,q,r\geq1,\ r=1}$} \inserted{$p\geq1,q\geq1,r=1$} *
(proof which moreover has the advantage
of \inserted{being expected to} apply without modification to the case of a
trace \unsure{$\tau$|$\sigma$|$\operatorname{Tr}$|$D$} in the sense of von Neumann - Dixmier, defined
on a self-adjoint algebra $\unsure{\mathcal A|\mathfrak A}$ \deleted{\unsure{strongly closed}}
\deleted{\unsure{uniformly closed}} uniformly closed \unsure{strongly}
(1) (\deleted{\unsure{the elements of}} \textsuperscript{(1)} the ``\deleted{\unsure{elements}} \inserted{operators} of power
$p$\textsuperscript{th} summable'' of $\unsure{\mathcal A|\mathfrak A}$ \inserted{-- in the sense of the trace $\unsure{\tau|\sigma|\operatorname{Tr}|D}$ --}
still being defined
as the set of $A\in\unsure{\mathcal A|\mathfrak A}$ such that
$\unsure{\tau|\sigma|\operatorname{Tr}|D}(|A|^p)<+\infty$.)\textsuperscript{\unsure{*|a|A}}.
\margin{?? It is moreover not difficult \deleted{\unsure{to}} then extend the same \deleted{\unsure{whole}} statement to the case where $r\geq1$}
(In the next number, we
shall moreover obtain the same result as
a consequence of a more general proposition).

We begin by proving that, if
$A_1,\ldots A_{2^n}$ are $2^n$ completely continuous operators,
($n$ an integer $\geq1$) then one has \textsuperscript{(1)}
\[
\|A_1 A_2\mathinner{\cdotp\cdotp\cdotp\cdotp} A_{2^n}\|_1\leq\|A_1\|_{2^n}\cdot\|A_2\|_{2^n}\ldots\|A_{2^n}\|_{2^n}\tag*{(11)}
\]
We prove it by induction on $n$\,;
this is true if $n=1$, \deleted{\unsure{one indeed has}} \inserted{(formula\quad)}, suppose
\deleted{\unsure{$\|A_1A_2\|_1\leq\|A_1\|_2$}}
it proved \inserted{\unsure{the|the} analogous formula} for $2^{n-1}$ factors, and let us prove
then (11). Writing $A_1\ldots A_{2^n}=(A_1A_2)\ldots(A_{2^n-1}A_{2^n})$
the left-hand side will, by the hypothesis
of induction, be bounded above by $\|A_1A_2\|_{2^{n-1}}\ldots\|A_{2^n-1}A_{2^n}\|_{2^{\unsure{n-1}}}$

\noindent\rule{\linewidth}{0.3pt}

(1) The reader will note that this result would suffice to
prove \deleted{\unsure{directly|exactly}} \inserted{directly}, as we shall do below,
formula \unsure{(6)|(6')} in the general case.
